% articulo-17-abstract.tex - Entorno abstract.
%
% Redefine el entorno abstract de la clase article para adaptarlo a la
% tipografía del documento. El formato por defecto de article (tamaño small,
% márgenes reducidos, título centrado) responde a las convenciones de los
% journals anglófonos: en un documento académico autónomo no hay razón para
% diferenciarlo visualmente del resto del paratexto preliminar.
%
% Formato resultante (Chicago § 1.73 adaptado):
%   · Cuerpo normalsize, interlineado según linestretch del documento.
%   · Título «Resumen» en negrita, alineado a la izquierda, en línea propia.
%   · Primer párrafo sin sangría de primera línea (\noindent + \ignorespaces).
%   · Párrafos subsiguientes con \parindent normal (heredado del documento).
%
% El \noindent inicial del entorno suprime la sangría solo del primer párrafo;
% los siguientes toman \parindent del contexto sin necesidad de intervención.
%
% Las palabras clave se imprimen en la maestra dentro del mismo entorno,
% inmediatamente tras el texto del abstract, como párrafo nuevo sin sangría.

\renewenvironment{abstract}{%
  \normalsize
  \noindent\textbf{Resumen}\par
  \noindent\ignorespaces
}{%
  \par
}
